Magnetic particles of various shapes and sizes are used in biomedical applications, robotics, nanomagnetism, electrotechnology, and education. Yet despite years of study, a theoretical description of the rolling motion of a solid magnetic sphere still eludes us. One of the reasons is that coupling magnetic fields with mechanical degrees of freedom creates a highly non-linear motion with non-holonomic constraints, which makes modeling these systems analytically or by simulations non-trivial. The Lagrange-d'Alembertian approach to holonomic systems remains one of the most versatile methods to prove central theorems and solve problems in analytical mechanics and electrodynamics. While non-holonomic systems are not accessible by the original theory, recent developments generalize the principle to certain dissipative systems. \newline Here, we present a successful application of these developments to a simple mechanical system comprising a magnetic solid sphere on an incline within an external magnetic field. Using numerical simulations, we characterized how the system depends on its initial parameters. Furthermore, our findings are in line with recent experiments, demonstrating that even very weak magnetic fields can cause unexpected motion of ferromagnetic objects. The simplicity of this example also makes it a promising candidate for undergraduate physics education where the treatment of non-holonomic variational principles is still widely uncommon.